% Revisione Ready
% Grammatica Ready
% Citazioni Ready

\chapter{Conclusioni e sviluppi futuri}\label{cha:conclusioni_e_sviluppi_futuri}

Il presente lavoro di tesi ha affrontato il problema combinatorio dell'Hitting Set in ipergrafi temporali, proponendo per esso una soluzione incrementale basata su una finestra temporale a scorrimento.
L'obiettivo del progetto è stato quello di evitare il ricalcolo completo ad ogni spostamento della finestra temporale, in modo da poter ottimizzare i tempi e la dimensione della soluzione stessa, introducendo un concetto di stabilità temporale per aumentare il riutilizzo dei nodi.
L'algoritmo sviluppato sfrutta il principio di stabilità temporale di una porzione delle interazioni per evitare il ricalcolo totale dell'Hitting Set.
Tale soluzione è stata implementata attraverso l'utilizzo della funzione \texttt{Window Slide Function} e mediante le relative strutture di supporto come \texttt{NodePresence} per tenere traccia dei nodi visitati in precedenza. 
La seguente funzione è stata successivamente valutata sperimentalmente attraverso l'utilizzo di quattro dataset reali strutturalmente diversi in modo da poter analizzare i trade-off strutturali e computazionali della soluzione.
L'utilizzo dell'approccio incrementale è risultato particolarmente efficace all'interno di scenari caratterizzati da distribuzioni temporali uniformi o più in generale senza concentrazioni troppo elevate, in tali contesti, l'algoritmo è risultato significativamente migliore rispetto ai classici approcci tradizionali, riducendo sia la dimensione media dell'Hitting Set che il numero totale di nodi distinti coinvolti nell'intera evoluzione della rete.
Di contro, sono stati però evidenziati dei limiti sistematici legati alla natura dei dati, in presenza di distribuzioni fortemente intermittenti e concentrate, il costo computazionale richiesto per la gestione delle strutture dati ausiliarie e per le fasi di potatura delle ridondanze non viene compensato da un reale guadagno nella qualità della soluzione.
Il ricalcolo della soluzione da zero, può risultare una soluzione praticabile all'interno di questi specifici scenari in quanto competitivamente preferibile in termini di tempo di esecuzione rispetto all'algoritmo sviluppato.
L'analisi ha infine confermato come l'ampiezza della finestra temporale $\epsilon$, il passo di scorrimento della finestra $\delta$ e il loro rapporto rappresentano ulteriori parametri fondamentali per riuscire a determinare l'efficacia della soluzione.

\section{Sviluppi futuri}

I limiti mostrati dai risultati ottenuti aprono la strada a diverse direzioni di ricerca futura:
\begin{itemize}
    \item \textbf{Parallelizzazione e Computazione Distribuita}: data l'attuale implementazione in modalità single-thread, un'estensione naturale prevede lo sviluppo dell'algoritmo per architetture parallele o contesti di stream processing, parallelizzando le fasi di verifica delle ridondanze e di aggiornamento dei nodi.
    \item \textbf{Strategie di Epurazione Adattive della Memoria}: per mitigare l'overhead della struttura NodePresence, si potrebbero studiare meccanismi di decadimento temporale o politiche di rimozione basate su un fattore di dimenticanza, eliminando dalla cronologia i nodi inattivi da troppo tempo.
    \item \textbf{Approssimazione Garantita e Rilassamento del Problema}: infine, risulterebbe molto interessante dal punto di vista teorico analizzare formalmente il fattore di approssimazione dell'algoritmo incrementale proposto rispetto all'ottimo globale e valutarne eventuali varianti euristiche e probabilistiche.
\end{itemize}